% Author notation from formal p23:3 footnote4 and p23:5.
The Minimum Circuit Size Problem with size parameter $s(n)$ ($\mathrm{MCSP}[s(n)]$) is defined as follows: given a string $x$ of length $N=2^n$, determine whether there exists an $n$-input $s(n)$ size circuit with $x$ as its truth table.
Formally, a probabilistic circuit deciding a language $L$ is a distribution $D$ over circuits that outputs the wrong answer with probability at most $1/\operatorname{poly}(n)$ on any input $x$, i.e., for all $x\in\{0,1\}^n$ it holds that $\Pr_{D\leftarrow D}[D(x)\ne L(x)]<1/\operatorname{poly}(n)$. We say a language $L$ is $s(n)$-sparse if for all sufficiently large $n$, $L\cap\{0,1\}^n$ has size at most $s(n)$.
Note that in contrast to robust classes like BPP, our result is actually sensitive to the error probability in the definition of probabilistic circuits, since the complexity overhead of error reduction is costly in the linear-size setting. For simplicity, we stick to the definition with error $1/\operatorname{poly}(n)$.
